\subsection{正则化}

命题 19. --- 设 G 为局部紧群，\(\beta\) 为 G 上一个 \(\neq 0\) 的相对不变正测度，B 为 \(e\) 在 G 中的邻域滤子的一个基，由紧邻域组成。对每个 V \(\in\) B，设 \(f_{V}\) 为 G 上的一个连续函数 \(\geqslant 0\)，其支集包含在 V 中，且满足 \(\int f_{V} d\beta = 1\)。若 \(\mu\) 是 G 上的一个测度，则在 \(\mathcal{M}(G)\)（赋有 \(\mathcal{K}(G)\) 中的紧收敛拓扑）中，

\[
\mu = \lim _ {\mathrm{V}} (\mu * f _ {\mathrm{V}}) \cdot \beta = \lim _ {\mathrm{V}} (f _ {\mathrm{V}} * \mu) \cdot \beta ,
\]

极限是沿 B 的截滤子取的。

对 \(\mathcal{C}(\mathrm{G})\) 中的紧收敛拓扑，\(f_{V} \cdot \beta\) 沿 B 的截滤子趋于 \(\varepsilon_{e}\)（\(\S2\), No.~7 引理 4 的推论 1）。因此 \(\mu = \lim_{\mathrm{V}} \mu * (f_{\mathrm{V}} \cdot \beta) = \lim_{\mathrm{V}} (f_{\mathrm{V}} \cdot \beta) * \mu\) 在 \(\mathcal{M}(\mathrm{G})\) 中成立，该空间赋有 \(\mathcal{K}(\mathrm{G})\) 中的紧收敛拓扑（\(\S3\), No.~3 命题 12 的推论）。

备注. --- 1) 由此我们看到，G 上的每个测度都是一些测度的极限，这些测度关于每个 Haar 测度都具有连续密度（对命题 19 中所指的拓扑，更不用说对淡拓扑了）。

\begin{enumerate}
\def\labelenumi{\arabic{enumi})}
\setcounter{enumi}{1}
\tightlist
\item
  若 G 可度量化，则 B 可取为一列邻域 \((\mathrm{V}_{n})\)。这时 \(\mu\) 是具有连续密度的测度序列 \((\mu * f_{\mathrm{V}_{n}}) \cdot \beta\) 的极限。*若 G 是一个实 Lie 群，则 \(f_{V_{n}}\) 可取为无穷次可微的；我们以后会看到，密度 \(\mu * f_{V_{n}}\) 这时是无穷次可微的。*
\end{enumerate}

命题 20. --- 我们保留命题 19 的假设与记号。设 \(p \in [1, +\infty[\) 与 \(g \in \mathrm{L}^p(\mathbf{G}, \beta)\)。则

\[
g = \lim _ {\mathrm{V}} g * ^ {\beta} f _ {\mathrm{V}} = \lim _ {\mathrm{V}} f _ {\mathrm{V}} * ^ {\beta} g
\]

在范数 \(N_{p}\) 的意义下，极限是沿 B 的截滤子取的。

只需应用命题 6 (iii) 以及 §2, No.~7 引理 4 的推论 3。

备注 3. --- 由命题 15，函数 \(g * f_{\mathrm{V}}\)、\(f_{\mathrm{V}} * g\) 属于 \(\overline{\mathcal{K}(\mathrm{G})}\)

推论. --- 设 W 为 \(\mathrm{L}^{1}(\mathrm{G},\beta)\) 的一个闭线性子空间。要使 W 是 \(\mathrm{L}^{1}(\mathrm{G},\beta)\) 的一个左（相应地，右）理想，必须且只需 W 在 G 的左（相应地，右）平移下不变。

假设 W 是一个左理想。设 \(s \in G\) 与 \(g \in W\)。我们有 \(\varepsilon_{s} * g = \lim_{\mathbb{V}} f_{\mathbb{V}} * (\varepsilon_{s} * g) = \lim_{\mathbb{V}} (f_{\mathbb{V}} * \varepsilon_{s}) * g\)，而 \((f_{\mathbb{V}} * \varepsilon_{s}) * g \in \mathbb{W}\)，因此 \(\varepsilon_{s} * g \in W\)，从而 \(\boldsymbol{\gamma}(s)g \in \mathbb{W}\)。反之，若 W 在左平移下不变，则 \(\mu *^{\beta}g \in W\)（对 \(\mu \in \mathcal{M}^{1}(\mathbf{G})\) 与 \(g \in W\)），因此 W 更是 \(\mathrm{L}^{1}(\mathrm{G}, \beta)\) 的一个左理想。对右理想可类似论证。

例. --- 我们取 \(\mathbf{G} = \mathbf{R}\)。我们用下式定义一个函数 \(\mathbf{F}_n \in \mathcal{K}(\mathbf{R})\)

\[
\begin{array}{l l} F _ {n} (x) = (1 - x ^ {2}) ^ {n} & \text { if } x \in [ - 1, 1 ] \\ F _ {n} (x) = 0 & \text { if } x \notin [ - 1, 1 ]. \end{array}
\]

设 \(\mathbf{A}_n = \int_{-1}^{+1}\mathbf{F}_n(x)dx\)，且 \(\mathbf{G}_n = \mathbf{A}_n^{-1}\mathbf{F}_n\)。立即可见测度 \(\mathbf{G}_n(x)dx\) 满足 §2, No.~7 引理 4 的推论 1 的条件。设 \(\mu\) 为 \(\mathbf{R}\) 上的一个测度，其支集包含在 \([-1/2, 1/2]\) 中。则

\[
\begin{array}{r l} & {(\mu * \mathrm{G} _ {n}) (x) = \int_ {\mathbf {R}} \mathrm{G} _ {n} (x - y) d \mu (y)} \\ & {\qquad = \mathrm{A} _ {n} ^ {- 1} \int_ {- 1 / 2} ^ {1 / 2} \mathrm{F} _ {n} (x - y) d \mu (y).} \end{array}
\]

若 \(-1 / 2\leqslant x\leqslant 1 / 2\)，则

\[
(\mu * \mathrm{G} _ {n}) (x) = \mathrm{A} _ {n} ^ {- 1} \int_ {- 1 / 2} ^ {1 / 2} [ 1 - (x - y) ^ {2} ] ^ {n} d \mu (y),
\]

因此 \(\mu * \mathbf{G}_n\) 在 \([-1/2, 1/2]\) 上与一个多项式重合。特别地，若 \(f\) 是一个连续函数，其支集包含在 \([-1/2, 1/2]\) 中，则 \(f * \mathbf{G}_n\) 在 \([-1/2, 1/2]\) 中与一个多项式重合；此外，由命题 5 (iv) 以及 §2, No.~7 引理 4 的推论 3，\(f * \mathbf{G}_n\) 一致收敛于 \(f\)。*若 \(f\) 是 \(\mathbf{C}^r\) 类的，则导数 \(\mathrm{D}^s(f * \mathbf{G}_n)\) 一致趋于 \(\mathrm{D}^s f\)（对 \(0 \leqslant s \leqslant r\)）。*

